\documentclass[a4paper]{arrowhead}

\usepackage[yyyymmdd]{datetime}
\usepackage{etoolbox}
\usepackage[utf8]{inputenc}
\usepackage{multirow}

\renewcommand{\dateseparator}{-}

\newcommand{\fparam}[1]{\textit{\textcolor{ArrowheadBlue}{#1}}}

%% Special references
\newcommand{\fref}[1]{{\textcolor{ArrowheadBlue}{\hyperref[sec:functions:#1]{#1}}}}
\newcommand{\mref}[1]{{\textcolor{ArrowheadPurple}{\hyperref[sec:model:#1]{#1}}}}
\newcommand{\pdef}[1]{{\textcolor{ArrowheadGrey}{#1 \label{sec:model:primitives:#1} \label{sec:model:primitives:#1s}}}}
\newcommand{\pref}[1]{{\textcolor{ArrowheadGrey}{\hyperref[sec:model:primitives:#1]{#1}}}}

\newrobustcmd\fsubsection[5]{
  \addtocounter{subsection}{1}
  \addcontentsline{toc}{subsection}{\protect\numberline{\thesubsection}function \textcolor{ArrowheadBlue}{#1}}
  \renewcommand*{\do}[1]{\rref{##1},\ }
  \subsection*{
    \thesubsection\quad
    #2 \textcolor{ArrowheadPurple}{#3} \\
    \small
    \hspace*{0.075\textwidth}\begin{minipage}{0.1\textwidth}
      \vspace*{1mm}
      Interface: \\
      \notblank{#4}{Input: \\}{}
      \notblank{#5}{Output: \\}{}
    \end{minipage}
    \begin{minipage}{0.825\textwidth}
      \vspace*{1mm}
      \textcolor{ArrowheadBlue}{#1} \\
      \notblank{#4}{\mref{#4} \\}{}
      \notblank{#5}{\mref{#5} \\}{}
    \end{minipage}
  }
  \label{sec:functions:#1}
}
\newrobustcmd\msubsection[2]{
  \addtocounter{subsection}{1}
  \addcontentsline{toc}{subsection}{\protect\numberline{\thesubsection}#1 \textcolor{ArrowheadPurple}{#2}}
  \subsection*{\thesubsection\quad#1 \textcolor{ArrowheadPurple}{#2}}
  \label{sec:model:#2} \label{sec:model:#2s}
}
%%

\begin{document}

%% Arrowhead Document Properties
\ArrowheadTitle{Event Subscribe HTTP/TLS/JSON}
\ArrowheadServiceID{event-subscribe}
\ArrowheadType{Interface Design Description}
\ArrowheadTypeShort{IDD}
\ArrowheadVersion{4.3.0}
\ArrowheadDate{\today}
\ArrowheadAuthor{Szvetlin Tanyi}
\ArrowheadStatus{RELEASE}
\ArrowheadContact{szvetlin@aitia.ai}
\ArrowheadFooter{\href{www.arrowhead.eu}{www.arrowhead.eu}}
\ArrowheadSetup
%%

%% Front Page
\begin{center}
  \vspace*{1cm}
  \huge{\arrowtitle}

  \vspace*{0.2cm}
  \LARGE{\arrowtype}
  \vspace*{1cm}

  \Large{Service ID: \textit{"\arrowid"}}
  \vspace*{\fill}

  % Front Page Image
  %\includegraphics{figures/TODO}

  \vspace*{1cm}
  \vspace*{\fill}

  % Front Page Abstract
  \begin{abstract}
    This document describes a HTTP/TLS/JSON variant of the Event Subscribe service.
  \end{abstract}

  \vspace*{1cm}

  \scriptsize
  \begin{tabularx}{\textwidth}{l X}
    \raisebox{-0.5\height}{\includegraphics[width=2cm]{figures/artemis_logo}} & {ARTEMIS Innovation Pilot Project: Arrowhead\newline
    THEME [SP1-JTI-ARTEMIS-2012-AIPP4 SP1-JTI-ARTEMIS-2012-AIPP6]\newline
    [Production and Energy System Automation Intelligent-Built environment and urban infrastructure for sustainable and friendly cities]}
  \end{tabularx}
  \vspace*{-0.2cm}
\end{center}
\newpage
%%

%% Table of Contents
\tableofcontents
\newpage
%%

\section{Overview}
\label{sec:overview}

This document describes the HTTP/TLS/JSON variant of the Event Subscribe Eclipse Arrowhead service, which enables subscribing to events.
Examples of this interaction is a consumer system would like to subscribe to an event.

This document exists as a complement to the \textit{Event Subscribe -- Service Description} document.
For further details about how this service is meant to be used, please consult that document.
The rest of this document describes how to realize the Event Subscribe service using HTTP \cite{fielding2014hypertext}, TLS \cite{rescorla2018transport} and JSON \cite{bray2014json}, both in terms of its functions (Section \ref{sec:functions}) and its information model (Section \ref{sec:model}).

\newpage

\section{Service Functions}
\label{sec:functions}

This section lists the functions that must be exposed by the Event Subscribe service in alphabetical order.
In particular, each subsection first names the HTTP method and path used to call the function, after which it names an abstract function from the Event Subscribe SD document, as well as input and output types.
All functions in this section respond with the HTTP status code \texttt{200 Created} if called successfully. The error codes are, \texttt{400 Bad Request} if request is malformed, \texttt{401 Unauthorized} if improper client side certificate is provided, \texttt{500 Internal Server Error} if Service Registry is unavailable.

\fsubsection{Query}{POST}{/eventhandler/subscribe}{SubscriptionRequest}{}

Called to subscribe to an event type, as exemplified in Listing \ref{lst:subscribe}.

\begin{lstlisting}[language=http,label={lst:subscribe},caption={An \fref{Event Subscribe} invocation.}]
POST /eventhandler/subscribe HTTP/1.1

{
  "eventType": "string",
  "filterMetaData": {
    "additionalProp1": "string",
    "additionalProp2": "string",
    "additionalProp3": "string"
  },
  "matchMetaData": true,
  "notifyUri": "string",
  "sources": [
    {
      "systemName": "string",
      "address": "string",
      "authenticationInfo": "string",
      "port": 0
    }
  ],
  "startDate": "string",
  "endDate": "string",
  "subscriberSystem": {
    "systemName": "string",
    "address": "string",
    "authenticationInfo": "string",
    "port": 0
  }
}

\end{lstlisting}
\newpage

\section{Information Model}
\label{sec:model}

Here, all data objects that can be part of the service calls associated with this service are listed in alphabetic order.
Note that each subsection, which describes one type of object, begins with the \textit{struct} keyword, which is meant to denote a JSON \pref{Object} that must contain certain fields, or names, with values conforming to explicitly named types.
As a complement to the primary types defined in this section, there is also a list of secondary types in Section \ref{sec:model:primitives}, which are used to represent things like hashes, identifiers and texts.

\msubsection{struct}{Metadata}

A JSON \pref{Object} which maps \pref{String} key-value pairs.

\msubsection{struct}{SubscriptionRequest}

This structure is used by a provider to broadcast its event.

\begin{table}[ht!]
\begin{tabularx}{\textwidth}{| p{5cm} | p{5cm} | X |} \hline
\rowcolor{gray!33} Object Field & Value Type      & Description \\ \hline
"eventType"        & \pref{String}                      & Type of the event. \\ \hline
"filterMetaData"   & \pref{Metadata}                    & Metadata filter.\\ \hline
"matchMetaData"    & \pref{Boolean}                     & Match by the provided metadata. \\ \hline
"notifyUri"        & \pref{String}                      & Uri to call by the Event Handler \\ \hline
"sources"          & \pref{Array}$<$\pref{System}$>$    & List of source systems \\ \hline
"startDate"        & \pref{DateTime}                    & Subscription start date \\ \hline
"endDate"          & \pref{DateTime}                    & Subscription end date \\ \hline
"subscriberSystem" & \pref{System}                      & Details of the subscriber system \\ \hline
\end{tabularx}
\end{table}

\msubsection{struct}{System}

This structure is used to describe a System \pref{Object}. 

\begin{table}[ht!]
\begin{tabularx}{\textwidth}{| p{5cm} | p{6cm} | X |} \hline
\rowcolor{gray!33} Object Field & Value Type                      & Description \\ \hline
"systemName"         & \pref{Name}     & System Name. \\ \hline
"address"            & \pref{String}   & Address. \\ \hline
"port"               & \pref{Port}     & Port. \\ \hline
"authenticationInfo" & \pref{String}   & Authentication Info. \\ \hline
\end{tabularx}
\end{table}

\subsection{Primitives}
\label{sec:model:primitives}

As all messages are encoded using the JSON format \cite{bray2014json}, the following primitive constructs, part of that standard, become available.
Note that the official standard is defined in terms of parsing rules, while this list only concerns syntactic information.
Furthermore, the \pref{Object} and \pref{Array} types are given optional generic type parameters, which are used in this document to signify when pair values or elements are expected to conform to certain types. 

\begin{table}[ht!]
\begin{tabularx}{\textwidth}{| p{3cm} | X |} \hline
\rowcolor{gray!33} JSON Type & Description \\ \hline
\pdef{DateTime}        & A moment in time.  \\ \hline
\pdef{Metadata}        & Metadata \\ \hline 
\pdef{String}             & A string identifier that is intended to be both human and machine-readable. \\ \hline
\pdef{System}           & An \pref{Object} describing a system. \\ \hline
\end{tabularx}
\end{table}

With these primitives now available, we proceed to define all the types specified in the Event Subscribe SD document without a direct equivalent among the JSON types.
Concretely, we define the Event Subscribe SD primitives either as \textit{aliases} or \textit{structs}.
An \textit{alias} is a renaming of an existing type, but with some further details about how it is intended to be used.
Structs are described in the beginning of the parent section.
The types are listed by name in alphabetical order.

\subsubsection{alias \pdef{DateTime} = \pref{String}}

Pinpoints a moment in time in the format of "YYYY-MM-DD HH:mm:ss", where "YYYY" denotes year (4 digits), "MM" denotes month starting from 01, "DD" denotes day starting from 01, "HH" denotes hour in the 24-hour format (00-23), "MM" denotes minute (00-59), "SS" denotes second (00-59). " " is used as separator between the date and the time.
An example of a valid date/time string is "2020-12-05 12:00:00"



\newpage

\bibliographystyle{IEEEtran}
\bibliography{bibliography}

\newpage

\section{Revision History}
\subsection{Amendments}

\noindent\begin{tabularx}{\textwidth}{| p{1cm} | p{3cm} | p{2cm} | X | p{4cm} |} \hline
\rowcolor{gray!33} No. & Date & Version & Subject of Amendments & Author \\ \hline

1 & 2020-12-05 & 1.0.0 & & Szvetlin Tanyi \\ \hline

\end{tabularx}

\subsection{Quality Assurance}

\noindent\begin{tabularx}{\textwidth}{| p{1cm} | p{3cm} | p{2cm} | X |} \hline
\rowcolor{gray!33} No. & Date & Version & Approved by \\ \hline

1 & & & \\ \hline

\end{tabularx}

\end{document}